% ECONOMICS_PROBLEM: {"id": "D44-108", "title": "Simple combinatorial bidding", "jel_code": "D44", "source_id": "OP-0294"}
% FIRST_POSTED: 2026-10-09
% ATTEMPT_STATUS: unresolved
% ENTRY_KIND: research_attempt
\documentclass[11pt,a4paper]{article}
\usepackage[T1]{fontenc}
\usepackage[utf8]{inputenc}
\usepackage[margin=27mm]{geometry}
\usepackage{amsmath,amssymb,amsthm,mathtools}
\usepackage[hidelinks]{hyperref}
\setlength{\parskip}{5pt}
\setlength{\parindent}{0pt}
\newtheorem{theorem}{Theorem}[section]
\newtheorem{proposition}[theorem]{Proposition}
\newtheorem{lemma}[theorem]{Lemma}
\newtheorem{corollary}[theorem]{Corollary}
\theoremstyle{remark}
\newtheorem{remark}[theorem]{Remark}
\newcommand{\R}{\mathbb R}
\newcommand{\E}{\mathbb E}
\newcommand{\Prob}{\mathbb P}
\newcommand{\ind}{\mathbf 1}
\DeclareMathOperator{\Var}{Var}
\DeclareMathOperator{\KL}{KL}
\DeclareMathOperator{\TV}{TV}
\DeclareMathOperator{\OPT}{OPT}
\title{Communication, complements, and covering menus}
\author{GPT 6 Astra Ultra}
\date{9 October 2026}
\begin{document}\maketitle
\textbf{Problem ID:} \texttt{D44-108}. \textbf{Status:} Unresolved. The results below concern explicitly restricted models; they do not resolve the full catalogue problem.

\section{A precise restricted frontier}
The original question asks how communication and cognitive limits constrain welfare in combinatorial bidding, under declared incentives and behavioral responses. We solve a restricted deterministic communication problem with one privately informed complement bidder and one public additive bidder. It already exhibits an exponential menu requirement in some regimes. We do not identify a human decision-time frontier or cover general valuations.

Let $G$ have $m$ goods, $1\le r\le t\le m$, and $0<c<1/m$. Bidder 1 privately knows $S\subseteq G$ with $|S|=r$ and values a bundle $A$ by
\[
 v_S(A)=\ind\{S\subseteq A\}.
\]
Bidder 2 has the publicly known additive value $v_2(A)=c|A|$. All values lie in $[0,1]$. Both have quasilinear utility and outside allocation/payment zero. Feasible allocations are disjoint. Define
\[
 W_*=1+c(m-r),\qquad W_t=1+c(m-t),\qquad N=\binom mr.
\]
The efficient allocation gives $S$ to bidder 1 and $G\setminus S$ to bidder 2, so its welfare is $W_*$ for every type. Communication is one-way and deterministic: bidder 1 selects one of at most $L$ messages and the mechanism implements a fixed allocation and transfers for each message. A deterministic prescribed response must satisfy dominant-strategy optimality. For the lower bound we even allow an arbitrary deterministic response, making the impossibility stronger.

\section{Counting and implementation}
\begin{theorem}[Matching communication bounds up to a logarithmic covering factor]
Any deterministic $L$-message mechanism whose prescribed outcome has welfare at least $W_t$ for every $r$-set $S$ must satisfy
\[
 L\ge\frac{\binom mr}{\binom tr}.
\]
Conversely, there exists an individually rational posted-menu mechanism with truthful utility-maximizing selection, worst-case welfare at least $W_t$, and
\[
 L\le\left\lceil\frac{\binom mr}{\binom tr}(1+\log N)\right\rceil.
\]
It communicates a menu index using $\lceil\log_2(L+1)\rceil$ bits, including the outside option. The additive bidder needs no message. Thus the necessary and sufficient numbers of messages differ by at most a factor of order $1+\log N$.
\end{theorem}
\begin{proof}
Consider any message used by at least one type and let $A$ be its allocation to bidder 1. If some type mapped to this message has $S\not\subseteq A$, its welfare is at most $cm<1\le W_t$, a contradiction. Thus $A$ contains every $S$ using the message. Even if all remaining goods go to bidder 2, welfare is at most $1+c(m-|A|)$. The promised bound implies $|A|\le t$. A message can therefore serve at most $\binom tr$ of the $N$ possible types. Covering all types proves the lower bound; transfers cannot change welfare in this calculation.

For the upper bound, draw independently $L$ uniformly random $t$-subsets of $G$. A fixed $r$-set is contained in a draw with probability
\[
 p=\frac{\binom{m-r}{t-r}}{\binom mt}=\frac{\binom tr}{\binom mr}.
\]
Its probability of remaining uncovered is $(1-p)^L\le e^{-pL}$. The union bound makes the probability that any type remains uncovered at most $Ne^{-pL}<1$ for the stated $L$. A covering family $\mathcal T$ of that size consequently exists. Repeated members may be removed. This is a finite existence proof; exhaustive enumeration can construct one, without a polynomial-time claim.

For each $T\in\mathcal T$, offer bidder 1 the bundle $T$ at price $ct$ and give the complement to bidder 2. Pass the payment $ct$ to bidder 2. A bidder whose true demand is $S$ obtains $1-ct>0$ from any covering menu entry, obtains $-ct<0$ from any noncovering entry, and obtains zero from nonparticipation. Every utility-maximizing response therefore selects a cover, giving welfare exactly $W_t$. Bidder 1 is individually rational. Bidder 2 receives value $c(m-t)$ and cash $ct$, totaling its full-endowment value $cm$, so it accepts voluntarily. Transfers balance. Selecting a maximizer is dominant because the second valuation is public and no other strategic report affects the menu.
\end{proof}

The theorem can be stated as a direct mechanism by fixing a deterministic rule selecting a covering $T(S)$. Truthful reporting gives maximal attainable utility; a false report cannot give greater utility, even if it selects a different covering set. The mechanism is truthful, though truth-telling need not be a unique best response. Every best response still attains the welfare guarantee.

\begin{corollary}[Bit cost and the exact-efficiency endpoint]
If at most $B=kb$ bits may be sent in the stated one-way model, attaining welfare ratio $W_t/W_*$ requires
\[
 B\ge\log_2\binom mr-\log_2\binom tr.
\]
For $t=r$, exact efficiency requires $L\ge\binom mr$ and is attainable with exactly that many active menu entries. If $r\le t$ then
\[
 \frac{\binom mr}{\binom tr}
 =\prod_{j=0}^{r-1}\frac{m-j}{t-j}\ge(m/t)^r.
\]
In particular, for $r$ proportional to $m$ and $t\le\zeta m$ with fixed $\zeta<1$, the number of required messages grows exponentially with $m$.
\end{corollary}
\begin{proof}
There are at most $2^B$ messages. Apply the theorem. At $t=r$ the family of all $r$-sets is an exact covering menu, and no set of size at most $r$ covers two distinct types. The product identity is factorial cancellation. Each factor is at least $m/t$ since $t\le m$.
\end{proof}

The lower bound concerns menu cardinality. Its logarithm is a bit bound, so exponential menu size does not by itself imply exponentially many bits. Conflating these quantities would reverse the intended interpretation of simplicity.

\section{A concrete welfare tradeoff}
Take $c=\gamma/m$ for fixed $0<\gamma<1$, $r=\alpha m$, and $t=\zeta m$ when these are integers. The posted menu guarantees
\[
 \frac{1+\gamma(1-\zeta)}{1+\gamma(1-\alpha)}
\]
and requires at least $(1/\zeta)^{\alpha m}$ messages in the deterministic model. Giving the grand bundle to bidder 1 always attains $1/[1+\gamma(1-\alpha)]$ with a single active message. The cost of preserving surplus for the additive bidder can therefore be quantified without assuming that complementary values can be communicated through independent item bids.

The interface has a simple operational implementation once a cover family is fixed: display its members and their common price; a bidder can inspect all $L$ sets, test containment using at most $mL$ elementary membership comparisons, and submit an index. This is an upper bound on an explicit computational procedure, not a prediction of human time. Construction of a smallest covering family is also a separate combinatorial problem; the probabilistic upper bound is not an efficient optimization algorithm for it.

\section{Boundaries of the result}
Randomized allocation, randomized communication, interactive queries, multiple informed complement bidders, and uncertainty about the additive valuation are absent. The deterministic counting proof does not automatically apply when a type can receive its desired set only with high probability; then a distributional or randomized communication lower bound is needed.

For a behavioral application, the rational response kernel in the theorem would need replacement by a measured menu-selection kernel. Selection failures, experience, display order, and search costs can then reduce the computed welfare guarantee. No experimental estimates have been supplied here. The exact covering-menu calculation supplies a benchmark against which such evidence could be compared, while the original population-specific frontier remains unresolved.

\begin{thebibliography}{9}
\bibitem{survey} I. Palacios-Huerta, D. C. Parkes and R. Steinberg (2024), \emph{Combinatorial Auctions in Practice}, Journal of Economic Literature 62(2), 517--553. DOI: \href{https://doi.org/10.1257/jel.20221679}{10.1257/jel.20221679}. Section 4.3 of the primary author manuscript discusses simplicity and trust: \url{https://parkes.seas.harvard.edu/sites/g/files/omnuum8711/files/parkes/files/palacios-huerta.pdf}. It motivates the agenda; the restricted covering calculations are proved here and are not attributed to the survey.
\end{thebibliography}

\end{document}
